% CVPR two-column report. Put the official CVPR author-kit files (cvpr.sty, ieeenat_fullname.bst)
% from https://github.com/cvpr-org/author-kit next to this file (or use the CVPR template on Overleaf).
% Without them the file still compiles with a plain two-column layout.
\documentclass[10pt,twocolumn,letterpaper]{article}
\IfFileExists{cvpr.sty}{\usepackage[final]{cvpr}}{\usepackage[margin=0.75in]{geometry}}
\usepackage{graphicx,booktabs,amsmath,amssymb,xcolor,url,adjustbox}
\usepackage[hidelinks]{hyperref}
\graphicspath{{../figures/}{figures/}}
\newcommand{\iffig}[3]{\IfFileExists{../figures/#1}{#2}{\IfFileExists{figures/#1}{#2}{#3}}}
% float placement: allow fuller pages so figures stay close to their text
\renewcommand{\topfraction}{0.9}\renewcommand{\bottomfraction}{0.8}\renewcommand{\textfraction}{0.07}
\renewcommand{\floatpagefraction}{0.8}\renewcommand{\dbltopfraction}{0.9}\renewcommand{\dblfloatpagefraction}{0.8}
\setcounter{topnumber}{3}\setcounter{bottomnumber}{2}\setcounter{totalnumber}{5}\setcounter{dbltopnumber}{2}
\newcommand{\todo}[1]{\textcolor{red}{[TODO: #1]}}
\newcommand{\tab}[1]{\IfFileExists{tables/#1.tex}{\input{tables/#1.tex}}{\todo{run make\_tables.py: #1}}}
\newcommand{\figw}[2]{\begin{figure*}[!tbp]\centering\iffig{#1}{\includegraphics[width=0.85\textwidth]{#1}}{\todo{figure #1}}\caption{#2}\label{fig:#1}\end{figure*}}
\newcommand{\fig}[3]{\begin{figure}[!tbp]\centering\iffig{#1}{\includegraphics[width=#2]{#1}}{\todo{figure #1}}\caption{#3}\label{fig:#1}\end{figure}}

\title{MLPs, CNNs, Autoencoders, SVMs and Clustering on CIFAR10}
\author{Bighnaraj Mohapatra \quad Roll no.\ M26CSA011\\
Repository: \url{https://github.com/bighnaraj-mohapatra/a3-cifar10} \quad Final commit: \texttt{c783b79}}

\begin{document}
\maketitle

\begin{abstract}
On CIFAR10 with one stratified split (seed 26011), a CNN and an MLP with matched parameter counts (816,202 vs 856,074) reach 87.54$\pm$0.45\% and 55.07$\pm$0.43\% test accuracy. All MLPs fit their training set ($\geq$97.95\%), so the MLP's deficit lies in generalisation, not in fitting. A fixed pixel permutation leaves the MLP unchanged (54.72$\pm$0.20\%) but drops the CNN to 52.05$\pm$0.42\%, so the CNN's advantage requires spatial locality; removing colour costs the CNN only 2.17 points. With 40\% wrong labels, an MLP reaches its best validation accuracy at epoch 5 while fitting only 11.87\% of the wrong labels, and memorises them later (98.78\% at epoch 200). For autoencoders, the best reconstruction does not give the most class information: a denoising AE with worse reconstruction gives the best linear probe (60.55\% vs 47.82\%). On the same AE features, an RBF SVM reaches 48.04\% test accuracy, while $k$-means clustering attains ARI 0.0506.
\end{abstract}

%=====================================================================
\section{Common setup}
\label{sec:setup}
From the 50,000 CIFAR10~\cite{krizhevsky2009cifar} training images we draw one stratified split of 45,000 training and 5,000
validation images (500 per class) with \texttt{train\_test\_split(stratify=y, random\_state=26011)};
26011 is the digit sequence of the roll number. The indices are saved in
\texttt{splits/split\_seed26011.npz} and every experiment loads that file. Validation data are used only
for model, architecture and hyper-parameter selection. The official test set is evaluated by a separate
script only after all choices are fixed.

Images are scaled to $[0,1]$ and standardised per channel with statistics of the (transformed) training
images. No data augmentation is used: random crops and flips encode a spatial prior, which would confound
the MLP--CNN comparisons of Q1 and Q2 and has no meaning for permuted images. All classifiers are trained
with AdamW (weight decay $5\cdot10^{-4}$), batch size 128 and a cosine schedule for 60 epochs. The
learning rate of every compared model is chosen from the same grid $\{3\cdot10^{-4},10^{-3},3\cdot10^{-3}\}$
by validation accuracy (seed 26011). For MLP-L and CNN-A the selected rate ($3\cdot10^{-3}$) lies at the upper edge of the grid, so larger rates were not explored. The checkpoint with the best validation accuracy is kept. Final
comparisons use seeds 26011--26013 and report mean $\pm$ sample standard deviation. Fairness criteria and
predictions were committed in \texttt{PREREGISTRATION.md} (commit \texttt{2931a26}) before the final runs.
\textbf{Reuse.} Components are introduced once and reused: MLP-B and CNN-A (Q1(a)) are the models of
Q2(a), and CNN-A is the fixed classifier of Q3(b). The Q3(a) autoencoders are the baseline of Q3(b), are
probed in Q3(c) and provide the Q4 representation (Table~\ref{tab:models} and the README reuse map).

%=====================================================================
\section{Q1: Why does an MLP struggle on CIFAR10?}
\subsection{Q1(a): separating the reason from an alternative}
\textbf{Chosen reason (H1).} The MLP has no spatial inductive bias. It sees an unordered 3,072-vector,
with neither locality nor weight sharing, so it must learn pixel relations from data and generalises
poorly even when it can fit the training set. \textbf{Alternative (H2).} The MLP under-fits because of
too little capacity or poor optimisation. The two hypotheses make different predictions. H2 predicts low
training accuracy, validation accuracy that rises with width, and large gains from better optimisation
(BN, tuned LR). H1 predicts high training accuracy with a large train--val gap, validation accuracy that
saturates as width grows, and a clear gain from a spatial prior at equal parameter count.
\textbf{Design.} (i) A width sweep MLP-S/B/L/XL (0.41M--4.21M parameters). (ii) An optimisation control,
MLP-B-plain (no BN, untuned LR). (iii) A prior-injection control, CNN-A, parameter-matched to MLP-B
(816,202 vs 856,074, within 5\%). All use the protocol of Sec.~\ref{sec:setup}. Compute is not matched
(Table~\ref{tab:models}); we report it rather than hide it.

\tab{tab_models}
\tab{tab_q1a}
\figw{q1a_curves.png}{Q1(a)/Q2(a): training and validation accuracy per epoch (seed 26011).}
\fig{q1a_width_sweep.png}{0.85\linewidth}{Q1(a): MLP width sweep vs the parameter-matched CNN-A.}

\textbf{Interpretation.} All MLPs fit the training set (maximum training accuracy 97.95\% for MLP-S and 100.00\% for the others), so H2 in its under-fitting form is not supported: the deficit is a generalisation gap (MLP-B: 100.00\% training vs 55.88\% best validation accuracy). Increasing the parameter count 10.2$\times$ (MLP-S to MLP-XL) raises test accuracy only from 53.61\% to 57.78\% (single seeds). Contrary to our pre-registered expectation of saturation, validation accuracy still rises slowly with width (54.24\% to 58.22\%): capacity matters, but its effect is small. Better optimisation (BN and a tuned LR) adds 1.54 points (MLP-B-plain 53.53\% vs MLP-B 55.07\%). In contrast, CNN-A, with 4.7\% fewer parameters than MLP-B, reaches 87.54$\pm$0.45\%, a 32.47-point gap, and even MLP-XL, with 5.2$\times$ CNN-A's parameters, is 29.76 points behind. The evidence favours H1. Because CNN-A uses 154$\times$ more MACs per image than MLP-B, this experiment alone cannot exclude compute as the cause; Q2(b) tests this.

\figw{q1b_memorisation.png}{Q1(b): training trajectory under 40\% label noise (left) and the
speed of fitting real vs random labels (right).}
\subsection{Q1(b): is structure learned before memorisation?}
\textbf{Setting.} MLP-Mem (3072--2048--2048--10, 10,510,346 parameters, no BN, no weight decay, no
dropout) is trained with Adam (constant LR $3\cdot10^{-4}$) for 200 epochs. A fixed 40\% of the
training labels (seed 26011) is replaced by a label drawn uniformly from the \emph{other} nine classes, so
fitting a corrupted label is unambiguously memorisation, as in random-label studies~\cite{zhang2017rethinking,arpit2017closer}. Validation labels stay clean. Every epoch we
measure accuracy on the clean-label subset, accuracy on the noisy subset against the assigned (wrong) and
the true labels, and validation accuracy. Two controls use 0\% and 100\% corrupted labels.
\textbf{Pre-registered rule.} Let $t^*$ be the epoch of peak validation accuracy. We conclude that
structure is learned first iff, at $t^*$, the noisy-subset fit to assigned labels is below 20\% and the
noisy-subset accuracy against true labels is above chance (10\%).

\tab{tab_q1b}

\textbf{Interpretation.} The pre-registered rule is met. At $t^*=5$ the network fits only 11.87\% of the corrupted labels (threshold 20\%), yet it predicts the \emph{true} class for 43.44\% of those corrupted images (chance 10\%); the clean-label subset is fitted at 58.88\%. The clean subset reaches 25/50/90\% fit at epochs 1/3/16 and the noisy subset at epochs 8/11/22. After $t^*$, memorisation takes over: by epoch 200 the noisy subset is fitted at 98.78\%, accuracy against its true labels falls to 0.29\%, and validation accuracy falls from 46.24\% to 35.84\% (test: 46.31\% at $t^*$ vs 35.22\% at the final epoch). Useful structure is therefore learned before memorisation, but the head start is short (6--8 epochs) and is measured for a single seed. The controls (Fig.~\ref{fig:q1b_memorisation.png}, right) show that the network has the capacity to memorise every label: with 100\% random labels it still approaches full training accuracy, but more slowly than with 0\% noise, where real structure is fitted first.

%=====================================================================
\section{Q2: What does the CNN use?}
\subsection{Q2(a): a fair comparison and a proposed reason}
We reuse MLP-B and CNN-A from Q1(a). The fairness criterion was fixed before the final results
(\texttt{PREREGISTRATION.md}): parameters within $\pm10\%$, identical data, preprocessing, optimiser,
epochs, LR grid and checkpoint rule, no augmentation, and three seeds. \textbf{Proposed reason:} CNN-A's
advantage comes from \emph{locality and weight sharing}. Its small filters apply the same local detectors
at every position and so exploit correlations between neighbouring pixels. Both models fit the training set (100.00\%), but CNN-A's best validation accuracy is 89.08\% against 55.88\% for MLP-B, so the difference lies in what each model generalises from, not in fitting capacity.

\subsection{Q2(b): same pixel values, disturbed spatial structure}
We apply one fixed random permutation of the 1,024 pixel positions (\texttt{default\_rng(26011)}),
identical for all channels and for every training, validation and test image, and retrain both models
with the same hyper-parameters. \textbf{Prediction stated beforehand.} MLP-B is unchanged within seed
noise: a fixed input permutation is equivalent to permuting its first-layer weights, whose i.i.d.\
initialisation is permutation-invariant. CNN-A drops substantially. If its advantage disappears, the
locality explanation of Q2(a) is supported. If it stays clearly ahead, it is not.

\tab{tab_q2}

\subsection{Q2(c): colour statistics or luminance structure?}
\textbf{Explanations.} E1: CNN-A relies mainly on colour statistics. E2: it relies mainly on luminance
spatial structure. \textbf{Experiment.} Both models are retrained and tested on grayscale images
($Y=0.299R+0.587G+0.114B$, replicated to three channels, so architecture and parameters are unchanged).
As a separately labelled diagnostic, the RGB-trained models are also evaluated on grayscale test images.
\textbf{Errors and confidence} (test set): per-class accuracy changes, new confusions, mean max-softmax for
correct and incorrect predictions, ECE with 15 bins~\cite{guo2017calibration}, and the share of errors made
with confidence $>0.9$ (\texttt{results/q2c\_confidence.json}).

\fig{q2c_reliability.png}{0.8\linewidth}{Q2(c): reliability diagrams for CNN-A (test set).}

\textbf{Conclusion linking (a)--(c).} Q2(a) attributed CNN-A's 32.47-point advantage to locality and weight sharing. Q2(b) tests this directly: a fixed pixel permutation keeps every pixel value and the CNN's compute, yet it leaves MLP-B unchanged within seed noise (55.07$\pm$0.43\% to 54.72$\pm$0.20\%), as predicted, while CNN-A drops 35.49 points to 52.05$\pm$0.42\%, below MLP-B. The advantage therefore needs the spatial arrangement of pixels, not extra parameters or compute. Q2(c) shows the CNN relies mainly on luminance structure rather than colour (E2): grayscale training costs CNN-A 2.17 points but MLP-B 8.38 points, and with RGB training on a grayscale test CNN-A still reaches 82.79$\pm$0.31\% against 40.64$\pm$0.26\% for MLP-B. Without spatial structure to use, the MLP leans more on colour. For CNN-A (seed 26011) the largest per-class drops under grayscale are cat ($-$4.8) and bird ($-$4.0), while the colour-typical classes change little (frog 0.0, ship +0.4), contrary to E1. Confidence does not track these errors: about 60\% of CNN-A's errors have confidence above 0.9 in all three conditions (0.600, 0.609, 0.591), and under the RGB-to-gray shift ECE rises from 0.0981 to 0.1337 while the mean confidence on wrong predictions barely changes (0.868 to 0.865).

%=====================================================================
\section{Q3: Autoencoders}
\subsection{Q3(a): bottleneck capacity}
The three fully convolutional AEs share every layer except the final $1\times1$ encoder and first
$1\times1$ decoder convolutions. The bottleneck is $z\in\mathbb{R}^{c\times4\times4}$ with
$c\in\{2,16,128\}$, i.e.\ $d\in\{32,256,2048\}$ (compression $96\times$, $12\times$, $1.5\times$). Training
is identical: per-pixel MSE on $[0,1]$, Adam with LR $10^{-3}$, cosine schedule, 50 epochs, batch
128, and the checkpoint with the lowest validation MSE. Selection uses validation MSE only.
\tab{tab_q3a}
\textbf{Choice.} Validation MSE decreases monotonically with bottleneck size (0.01231, 0.00212 and 0.00037 for $d=32$, 256 and 2048), so the reconstruction-only criterion selects $d=2048$. Its test MSE is 0.00036 (PSNR 35.32\,dB), matching its validation MSE. The choice should be read with care: $d=2048$ compresses only 1.5$\times$, close to an identity map, so reconstruction alone rewards the least compression. Training and validation MSE stay together for all three models (Fig.~\ref{fig:q3a_mse_curves.png}; for $d=2048$ the validation MSE is even slightly below the running training MSE, which is averaged during each epoch), so there is no sign of over-fitting.
\fig{q3a_mse_curves.png}{0.85\linewidth}{Q3(a): reconstruction MSE per epoch.}

\subsection{Q3(b): denoising}
\textbf{Corruption.} Additive Gaussian noise, $\sigma=0.2$ on $[0,1]$, clipped. It is a corruption a
reconstruction model can plausibly undo, because i.i.d.\ noise lies off the image
manifold~\cite{vincent2008denoising}. $\sigma$ was fixed in advance. A validation scan
(\texttt{results/q3b\_sigma\_scan\_val.json}) confirms that it hurts the classifier. \textbf{Protocol.}
The DAE has the same architecture as the selected clean AE and is trained on noisy inputs with clean
targets, without labels. Fresh noise is drawn each batch. Validation and test use fixed seeded noise.
The fixed classifier is the clean CNN-A (Q2(a), seed 26011). AE(clean) is a control for accuracy lost to
reconstruction blur alone.
\tab{tab_q3b}
\figw{q3b_examples.png}{Q3(b): test images: clean, noisy, AE(noisy), DAE(noisy).}
\textbf{Decision.} $\sigma=0.2$ was fixed in advance; on validation it reduces CNN-A from 89.08\% to 16.62\% (70.54\% at $\sigma=0.05$, 38.80\% at $\sigma=0.1$), so it removes most of what the classifier uses. On test, the DAE raises CNN-A from 17.11\% (noisy) to 62.11\%, recovering 45.00 of the 70.06 points lost, and agrees with the clean prediction on 64.34\% of the images; the clean AE recovers much less (32.03\%). Denoising therefore recovers a substantial part, but not all, of the information the classifier needs. Part of the remaining gap need not be lost information: the DAE also lowers accuracy on \emph{clean} inputs to 52.35\% (vs 85.45\% for the clean AE), which shows that CNN-A, trained on sharp images, is mismatched to the DAE's smoothed outputs. With a fixed classifier, these two causes cannot be separated.

\subsection{Q3(c): class information vs reconstruction}
Encoders are frozen. For each representation (the three AEs, the DAE, an untrained encoder of the selected
architecture, and raw pixels) we fit a multinomial logistic-regression probe (L2 strength chosen on
validation) and a kNN probe ($k$ chosen on validation) on the 45k training features, then report test
accuracy. Probes after PCA to 32 dimensions control for the different dimensionalities.
\tab{tab_q3c}
\textbf{Do the two quantities agree?} No. Among the three clean AEs the linear probe follows reconstruction (33.46\%, 40.58\%, 47.82\% as test MSE falls from 0.01232 to 0.00036), but this ordering is largely a dimension effect: after PCA to 32 dimensions the range shrinks to 33.50--37.86\%, and kNN reverses it (39.75\%, 30.18\%, 24.08\%). Two comparisons break the link outright. The DAE reconstructs clean images worse than AE-2048 (MSE 0.00287 vs 0.00036) yet gives the best linear (60.55\%) and PCA-32 (43.04\%) probes. An untrained encoder of the same architecture reaches 50.10\%, above every clean AE. Reconstruction quality is therefore not a reliable proxy for class information here. Each entry is a single run.

%=====================================================================
\section{Q4 (bonus): SVMs and clustering}
\textbf{Fixed representation.} The $d=256$ bottleneck of the Q3(a) AE, standardised, fixed in \texttt{PREREGISTRATION.md} before Q3 was run. We keep it rather than the reconstruction-selected $d=2048$, which is close to an identity map (Q3(a)). It was learned
without labels, so training-set features carry no label memorisation. \textbf{Fixed subset.} 10,000
training images (1,000 per class, seed 26011), chosen before any model selection and used for every SVM
and clustering fit. Validation uses the full 5k and test the full 10k.

\subsection{Q4(a): $C$ and kernel}
Kernels: linear, RBF and polynomial (degree 3, $c_0=1$), with $C\in\{0.01,0.1,1,10,100\}$ and $\gamma$
fixed at the numeric value of \texttt{'scale'}. Selection uses validation accuracy only; the selected SVM is
evaluated once on test.
\tab{tab_q4a}
\figw{q4a_sweep.png}{Q4(a): validation/training accuracy and support-vector fraction vs $C$.}
\textbf{Under- and over-fitting.} Small $C$ under-fits for every kernel (RBF, $C=0.01$: 33.90\% train, 31.36\% validation; every point is a support vector). Increasing $C$ raises training accuracy to 100.00\% for the RBF and polynomial kernels while validation accuracy saturates: RBF peaks at 48.06\% ($C=10$) and stays at 47.90\% at $C=100$, and the polynomial kernel over-fits visibly (47.20\% at $C=1$ falling to 44.80\% at $C=100$). The support-vector fraction stays between 0.928 and 1.000 for all RBF and polynomial settings: nearly every training point lies on or inside the margin, so the classes overlap strongly in this representation, and a larger $C$ buys training accuracy by bending the boundary around individual points rather than by widening the margin. The linear kernel did not converge within $2\times10^6$ iterations for $C\geq1$; those rows (including 28.14\% validation at $C=100$) are optimisation artefacts and are not interpreted. The selected SVM (RBF, $C=10$) reaches 48.04\% on the test set.

\subsection{Q4(b): difficult pairs}
The three pairs with the largest symmetric test confusion are analysed with the pairwise SVM (same kernel,
$C$ and numeric $\gamma$; exactly the one-vs-one classifier of the selected SVM).
\tab{fig_q4b}
The three pairs with the most symmetric confusions are cat--dog (410), automobile--truck (381) and bird--deer (311). In each, almost every training point is a support vector (SV fraction 0.9585, 0.9215, 0.877), but few are bounded (9, 0, 63): with $C=10$ the RBF boundary passes through nearly every training point exactly at the margin. Most test points fall inside the margin band $|f|<1$ (94.25\%, 91.40\%, 84.45\%), and so do the errors (96.7\%, 99.3\%, 93.0\%): they are low-margin rather than confident mistakes. The geometry explains this: class centroids are 1.53, 2.72 and 1.88 apart, while each class spreads about 14--17 units around its centroid (Fisher ratios 0.0044, 0.0123, 0.0087). The pairwise test accuracies (61.65\%, 73.20\%, 67.90\%) order the same way as the Fisher ratios. In Fig.~\ref{fig:q4b}, the bounded support vectors are low-contrast images or show a small object against a large uniform background (e.g.\ a deer in a misty field, a bird's neck against grass). The test images nearest the boundary share background, colour and framing across the two classes: close-up furry faces for cat--dog, front or rear views of boxy vehicles for automobile--truck, and brown animals on natural backgrounds for bird--deer. This is consistent with a representation trained for reconstruction, which encodes global colour and layout more than object identity.

\subsection{Q4(c): clustering vs classification}
$k$-means and diagonal GMMs with $k\in\{5,8,10,12,15,20\}$ are fitted on the subset without labels. The
method and $k$ are chosen by validation silhouette. Test representations are assigned without labels, and
labels are used only to evaluate.
\tab{tab_q4c}
\fig{q4c_pairs.png}{0.75\linewidth}{Q4(c): per class pair, SVM confusion rate vs clustering ARI (test).}
\textbf{Disagreement.} The label-free rule selected $k$-means with $k=5$. On test, the clusters agree only weakly with the classes (ARI 0.0506, NMI 0.0802), while the SVM on the same features reaches 48.04\% (and $k=5<10$ caps the possible agreement). The clearest disagreement is ship--truck: the SVM rarely confuses the two (pair confusion rate 0.0665), but the clustering does not separate them at all (pair ARI 0.0083). In Table~\ref{tab:q4c}, 618 ships and 503 trucks fall in the same cluster c4, which is the largest cluster for all four vehicle classes (airplane 401, automobile 318) and for no animal class. Automobile--dog shows the same pattern (SVM 0.0175 vs ARI 0.0429). Clustering follows the dominant variation in the representation, which groups the vehicles; the SVM uses labels to find the directions that separate ship from truck inside that group.

{\small
\IfFileExists{ieeenat_fullname.bst}{\bibliographystyle{ieeenat_fullname}}{\bibliographystyle{plain}}
\bibliography{references}
}
\end{document}
